% Source: https://artofproblemsolving.com/community/c3h1787416_extension_question
% Stable identity: aops-topic-1787416-q01
Triangle $ADC$ has a right angle at $D$, so
\[
AD^2+DC^2=13^2=169,\qquad \dfrac{1}{2}\,AD\cdot DC=30\ \Longrightarrow\ AD\cdot DC=60 .
\]
Then $(AD+DC)^2=169+120=289$ and $(AD-DC)^2=169-120=49$, so $AD+DC=17$, $|AD-DC|=7$, and $\{AD,DC\}=\{12,5\}$.

\textbf{Case $AD=12$, $DC=5$.} In right triangle $ADB$, $BD=\sqrt{15^2-12^2}=9$. In an acute triangle the foot of an altitude lies between the other two vertices, so $BC=BD+DC=14$ and
\[
[ABC]=\dfrac{1}{2}\cdot14\cdot12=84 .
\]

\begin{asy}
usepackage("amsmath");
size(7cm);
defaultpen(fontsize(10pt));
pair B=(0,0),C=(14,0),A=(9,12),D=(9,0);
draw(A--B--C--cycle,linewidth(0.9));
draw(A--D,linewidth(0.8));
draw((9,0)--(9,0.7)--(8.3,0.7)--(8.3,0),linewidth(0.5));
fill(A--D--C--cycle,lightgray+opacity(0.5));
draw(A--B--C--cycle,linewidth(0.9));draw(A--D,linewidth(0.8));
label("$A$",A,N);label("$B$",B,SW);label("$C$",C,SE);label("$D$",D,S);
label("$15$",(A+B)/2,NW);label("$13$",(A+C)/2,NE);
label("$12$",(9,6),W);
label("$9$",(4.5,0),S);label("$5$",(11.5,0),S);
label("\mbox{}",(min(currentpicture,user=true)+max(currentpicture,user=true))/2,black); // keeps the text after the figure black when LaTeX typesets the labels inline
\end{asy}

\textbf{Case $AD=5$, $DC=12$.} Then $BD=\sqrt{225-25}=10\sqrt2$, and the area $\dfrac{1}{2}\cdot5\,(10\sqrt2+12)=30+25\sqrt2$ is not an integer. (This triangle, with $BC=10\sqrt2+12$, is not even acute: $15^2+13^2=394<BC^2$.)

Hence the area of triangle $ABC$ is $\boxed{84}$; it is the $13$-$14$-$15$ triangle.
